\subsection{The prime losses at the selected translates}
\label{orb:selected-loss}

Write \(B=n_Tmb\), where \(a_E=xb\). At an actual point let
\(\sigma_p=\operatorname{Sig}(p)\), evaluated with its recorded integer
weights \(a,b_{\rm top}\), and put \(\theta=x^{-\varepsilon}\).
The candidates are the primes in \code{sAt}. The selected weights satisfy
\begin{equation}\label{orb:selected-W}
 \sum_{p\ \mathrm{candidate}}\frac{\log p}{p}
       \bigl(2{\bf1}_{\sigma_p>\theta}+4\sigma_p\bigr)
       \le x^{-9}.
\end{equation}
This is the predicate \code{ChW\_at}. On a kit carrying \code{WOK},
\code{wreq\_of\_wok} identifies the arithmetic share with the point
share for each \code{Pt.From} with canonical representatives, valuations
and zero weights off the tops. The equality uses both weight vectors.
The same centres therefore give \eqref{orb:selected-W} at all the
translates supplied by Proposition~\ref{scope:reselected-hosts}.
More explicitly, fix the complete windows \(S\) and select the entire
host \(H_s\), including both sets, at each scale. In the staged range
of the original point and weight existence theorems the order is
\[
   \exists x_1\quad\forall s\in S,\ s.x\ge x_1\quad
   \exists(a_s,b_s)\ \text{on }H_s\quad
   \forall\,0\le u<n_s\quad\exists\,p_u\ \text{from }H_s,a_s,b_s .
\]
First build the temporary family with provisional in-window centres.
At a kit point satisfying \code{WeightVectorExists}, the theorem
\code{weightWindowRequirement\_host\_of\_point} supplies duty
centres on its prescribed \(H_s\). Apply \code{WClosed} to this family.
Its windows are exactly \(S\), so an index at each late \(s\in S\)
recovers the duty centres at that same scale. Apply \code{kit\_of\_points}
again with these centres to obtain all the \(p_u\).
In particular selecting \(H_s\) by Proposition~\ref{orb:selection}
preserves its overlap and large-order properties throughout both
applications. The scale bound may depend on the prescribed sequence
of hosts.
\src{src/Development/AbelianVarietyTheory/Proofs/GlobalInterpolation/WeightWindowRequirementHostOfPoint.lean}{17}{returning from a point to its prescribed host}
\src{src/Development/AbelianVarietyTheory/Proofs/GlobalInterpolation/WeightRequirementOfValid.lean}{32}{the share inequality at a point of the common kit}.

\begin{proposition}[The signed prime loss on a selected point]\label{orb:selected}
Fix an original family \code{FamP} at exponent \(\varepsilon\).
There are \(C\ge0\) and \(x_0\) such that at every \(x\ge x_0\) and
every family point with \(n_T>0\), \eqref{orb:selected-W}, and
\[
  \sum_{p\ \mathrm{large\ order}}\frac{\log p}{p}
       \#\{\text{bad tops for }p\}\le n_Tx^{-2+\varepsilon},
\]
each tied reading of the asks with its required finite cutoff satisfies
\[
          \mathcal A_u-G_u\le CB,\qquad I_u\le Bx^{-8}.
\]
In particular, for a common kit whose points satisfy these conditions,
every \(\kappa\ge0\) gives
\[
   \mathbb E_{u\in A_0}(\mathcal A_u-G_u+\kappa I_u)
        \le (C+\kappa x^{-8})B .
\]
\end{proposition}
\begin{proof}
For an orbit representative \(v\) let \(N_v\) be its number of tops.
The ordinary parameter choice gives
\(S_v\le 2mN_v a_E/p\) by \code{svBound\_of}. The defining
estimate for the orbit term now gives
\(\mathrm{io}_p\le 4m n_T a_E\,\sigma_p/p\).
All shares are nonnegative by the definition of the local infimum.
Multiply by \(\log p\), sum and use \eqref{orb:selected-W}.
The identification \code{IO\_at} gives \(I_u\le
m n_T a_E x^{-9}=Bx^{-8}\).

For the gain comparison set \(b_o=b^{x^{-\varepsilon/2}}\).
For all sufficiently large \(x\), all the family indices, all
in-window vectors and every inert prime \(p\le X''\), the local
computed functions satisfy
\[
 \sum_j\sum_{t\in T_p}
    \bigl(\mathrm{askP}_{jtp}-\mathrm{lo}_p\bigr)-m\,\mathrm{fl}_p
 \le \sum_j\mathrm{gainP}_{jp}-\mathrm{\tau P}_p .
\]
Here the constants in \(\mathrm{lo}_p,\mathrm{fl}_p\) are uniform and
the bound includes the recorded vectors. To see the required inputs,
\code{family\_first\_lemma} supplies the count at good representatives,
and \code{family\_sizes\_bound} supplies the numerical prime windows.
Its positive coefficient \(K_{\rm La}\) follows at any late point
from \(\mathcal L_a=K_{\rm La}N_1^2\log a_E>0\).
The numerical estimates \code{numStar\_ev}, \code{facts\_ev} and
the representative count give \code{starAll\_of}. The additional
exponent in \code{gain\_bound\_at\_of\_star} is bounded directly
by rounding its recorded value
\(\lceil C_e f k^2\log_p a_E\rceil\).
For \(p\le X''\le a_E\) one has \(fk^2\log_p a_E\ge1\), so
replacing \(C_e\) by \(C_e+1\) absorbs the rounding.

On a served prime the point's fill identities identify these asks
and gains with the finite ledger. Enlarging the constant \(c\)
in its floor terms to \(4|c|+1\) gives the displayed inequality
with the ledger floors. Missing tops on served primes enter the
large-order bad-top sum. On a first-kind prime \code{firstHyp\_at}
puts the prime in the small-order set or in one of the specified
bands about \(X'\). For a second-kind prime
\(\sigma_p>\theta\), and its entire cost is at most
\[
    4m n_T a_E\sum_{\mathrm{second}\ p}\frac{\log p}{p}
        \le 2Bx^{-8}.
\]
The original numerical discriminant estimate
\code{joint\_discriminant\_eventually} combines these bounds with
the high-prime and floor terms, giving
\(\mathcal A_u-G_u\le(3C_{\rm par}+13)B\).
The parameter windows bound \(C_{\rm par}\) uniformly in \(x,u\).
Taking the uniform mean proves the last assertion. The independent
Lean implementation is \code{OrbitLossAtSelectedPoints.lean}.
\end{proof}

For the original construction, the large-prime step in
\eqref{orb:selected-W} involves a further finite average.
When \(p\ge p_0(k)\), \code{kapI} is the infimum of the losses
that satisfy \code{LPprime} in every admissible local DVR setting
and for every dominant lift of the prescribed top weights. In
this predicate the threshold is \(\lceil t_{\rm rec}\rceil\), the
allowed error level is \(D_{\rm I}\), and the accompanying level
is \(8(D_{\rm I}+1)\). On an orbit of length \(n'\) with \(N\) tops,
let \(d_y\) be the backward distance from a non-type-V top \(y\)
to the nearest type-V top, with value zero if none exists.
With cyclic indices the threshold \code{tIf} is
\[
\begin{aligned}
 t_{\rm rec}={}&
   \sum_{\substack{y\ {\rm non\text{-}V}\\ d_y>0}}
        \min\{b_y,a_{y-d_y}/p^{d_y}\}\\
 &+K_1NN_1\bigl(a_E/p^2+\mathcal L_a a_E/(pk)\bigr)\\
 &+K_2N\bigl(Q_pc_6+kQ_p+k\log_p a_E+c_G\bigr)+2n'c_G ,
 \qquad Q_p=p^{G'} .
\end{aligned}
\]
Here \(c_6\) is the local \code{c6f} at this prime and orbit length.
Thus each contributing top reads one predecessor and its actual
integer weights. The segment theorem has its own real threshold
\(t_{\rm seg}\) and holds for every integer \(t\ge t_{\rm seg}\).
The inequality \code{tI\_le\_tIf} gives
\(t_{\rm seg}\le t_{\rm rec}\le\lceil t_{\rm rec}\rceil\).
The original
\code{type\_one\_hypothesis\_of\_segments} transfers the local
segment estimate to the infimum using this inequality.
For each fixed good \(v\), the resulting bound over its local
weight box is
\[
 \mathbb E_{a,b_{\rm top}}\mathrm{kapI}_{p,v}
      \le C_m N_v N_1 x^{-13}\le N_v x^{-12}.
\]
The disjoint top fibres have \(\sum_vN_v\le n_T\).
The finite box argument of \code{meanEx\_of} followed by
\code{weight\_condition\_of\_mean\_excess} selects one common pair
of vectors satisfying \eqref{orb:selected-W}. This is where the
original construction uses its local representation estimates.
The independent bounds above apply once that pair is realized.
For an arbitrary prescribed vector the definition alone gives
\(0\le\mathrm{kapI}\le1\).
In particular the elementary construction by rounding the two
endpoints only establishes the window inequalities. The local
\code{LPprime} estimate over the weight box is an input from the
original proof in the duty-centre construction.

The selection of Proposition~\ref{orb:selection} makes the bad-top
cost zero. On the common selected kit, \(B\le n_T a_E\le P_0\) in
the notation of Corollary~\ref{scope:one-word}. Thus its two prime
loss requirements already hold at every eligible translate, before
estimating any word height. For
Corollary~\ref{trmean:loss} one can take \(e_{\rm loss}=0\) and
\(c_{\rm loss}=\max(C,1)\). The remaining comparison in
\eqref{scope:average} concerns \(G_u-\mathcal M_u\) at the
recorded word and target lattices.
